% GENERATED FILE. Edit canonical lessons, then run npm run generate:books.
% canonical-source-hash: fnv1a64:94947cf6ecb86dc0
% canonical-lessons: ML-C231-caruka, ML-C231-muttuka, ML-C231-kuniyuka, ML-C231-parikkuka, ML-C231-karakkuka, ML-R231-first-pass-words-from-chapters-223-226, ML-R231-second-pass-words-from-chapters-227-231

\chapter{To lean, To knock, To bend down, To pluck, To milk}
\label{ch:ml-to-lean-to-knock-to-bend-down-to-pluck-to-milk}
\hlchaptermodality{231}

\begin{chapteropening}
\textbf{By the end of this chapter:} \emph{I can say to lean, to knock, to bend down, to pluck and to milk.}

Complete the last lesson of chapter 231: I can say to lean, to knock, to bend down, to pluck and to milk.
\end{chapteropening}

\section[\texorpdfstring{\ml{ചാരുക} (cāruka)}{ചാരുക (cāruka)}]{\ml{ചാരുക} (cāruka) — to lean (against)}

\label{lesson:ML-C231-caruka}



\begin{quote}
Before the new one: say the Malayalam for to solve, then the Malayalam for to be bored.
\end{quote}

\subsection*{You'll want to know: \ml{ചാരുക}}
\textbf{\ml{ചാരുക}} — \emph{cāruka} — \textquotedblleft{}to lean (against)\textquotedblright{}.

\begin{grammarlens}[title={What you've built}]
One more word for this chapter.
\end{grammarlens}

\subsection*{Your turn}
\begin{itemize}
\raggedright
  \item \emph{Say it:} \emph{cāruka}
  \item \emph{Say it:} \emph{cāruka}, once more
  \item \emph{Say it:} say \emph{cāruka}
  \item \emph{Recall:} say the Malayalam for to solve, then the Malayalam for to be bored, then say \emph{cāruka} again
\end{itemize}

\begin{tcolorbox}[breakable,colback=teal!4,colframe=teal!35!black,title={Before you move on}]
What does \ml{ചാരുക} mean? (\textquotedblleft{}To lean (against)\textquotedblright{}.) Say it once more.
\end{tcolorbox}
\section[\texorpdfstring{\ml{മുട്ടുക} (muṭṭuka)}{മുട്ടുക (muṭṭuka)}]{\ml{മുട്ടുക} (muṭṭuka) — to knock (on a door)}

\label{lesson:ML-C231-muttuka}



\begin{quote}
Before the new one: say the Malayalam for to be bored, then the Malayalam for to lean.
\end{quote}

\subsection*{You'll want to know: \ml{മുട്ടുക}}
\textbf{\ml{മുട്ടുക}} — \emph{muṭṭuka} — \textquotedblleft{}to knock (on a door)\textquotedblright{}.

\begin{grammarlens}[title={What you've built}]
One more word for this chapter.
\end{grammarlens}

\subsection*{Your turn}
\begin{itemize}
\raggedright
  \item \emph{Say it:} \emph{muṭṭuka}
  \item \emph{Say it:} \emph{muṭṭuka}, once more
  \item \emph{Say it:} say \emph{muṭṭuka}
  \item \emph{Recall:} say the Malayalam for to be bored, then the Malayalam for to lean, then say \emph{muṭṭuka} again
\end{itemize}

\begin{tcolorbox}[breakable,colback=teal!4,colframe=teal!35!black,title={Before you move on}]
What does \ml{മുട്ടുക} mean? (\textquotedblleft{}To knock (on a door)\textquotedblright{}.) Say it once more.
\end{tcolorbox}
\section[\texorpdfstring{\ml{കുനിയുക} (kuniyuka)}{കുനിയുക (kuniyuka)}]{\ml{കുനിയുക} (kuniyuka) — to bend down, to stoop}

\label{lesson:ML-C231-kuniyuka}



\begin{quote}
Before the new one: say the Malayalam for to lean, then the Malayalam for to knock.
\end{quote}

\subsection*{You'll want to know: \ml{കുനിയുക}}
\textbf{\ml{കുനിയുക}} — \emph{kuniyuka} — \textquotedblleft{}to bend down, to stoop\textquotedblright{}.

\begin{grammarlens}[title={What you've built}]
One more word for this chapter.
\end{grammarlens}

\subsection*{Your turn}
\begin{itemize}
\raggedright
  \item \emph{Say it:} \emph{kuniyuka}
  \item \emph{Say it:} \emph{kuniyuka}, once more
  \item \emph{Say it:} say \emph{kuniyuka}
  \item \emph{Recall:} say the Malayalam for to lean, then the Malayalam for to knock, then say \emph{kuniyuka} again
\end{itemize}

\begin{tcolorbox}[breakable,colback=teal!4,colframe=teal!35!black,title={Before you move on}]
What does \ml{കുനിയുക} mean? (\textquotedblleft{}To bend down, to stoop\textquotedblright{}.) Say it once more.
\end{tcolorbox}
\section[\texorpdfstring{\ml{പറിക്കുക} (paṟikkuka)}{പറിക്കുക (paṟikkuka)}]{\ml{പറിക്കുക} (paṟikkuka) — to pluck, to pick (fruit, flowers)}

\label{lesson:ML-C231-parikkuka}



\begin{quote}
Before the new one: say the Malayalam for to knock, then the Malayalam for to bend down.
\end{quote}

\subsection*{You'll want to know: \ml{പറിക്കുക}}
\textbf{\ml{പറിക്കുക}} — \emph{paṟikkuka} — \textquotedblleft{}to pluck, to pick (fruit, flowers)\textquotedblright{}.

\begin{grammarlens}[title={What you've built}]
One more word for this chapter.
\end{grammarlens}

\subsection*{Your turn}
\begin{itemize}
\raggedright
  \item \emph{Say it:} \emph{paṟikkuka}
  \item \emph{Say it:} \emph{paṟikkuka}, once more
  \item \emph{Say it:} say \emph{paṟikkuka}
  \item \emph{Recall:} say the Malayalam for to knock, then the Malayalam for to bend down, then say \emph{paṟikkuka} again
\end{itemize}

\begin{tcolorbox}[breakable,colback=teal!4,colframe=teal!35!black,title={Before you move on}]
What does \ml{പറിക്കുക} mean? (\textquotedblleft{}To pluck, to pick (fruit, flowers)\textquotedblright{}.) Say it once more.
\end{tcolorbox}
\section[\texorpdfstring{\ml{കറക്കുക} (kaṟakkuka)}{കറക്കുക (kaṟakkuka)}]{\ml{കറക്കുക} (kaṟakkuka) — to milk (a cow)}

\label{lesson:ML-C231-karakkuka}



\begin{quote}
Before the new one: say the Malayalam for to bend down, then the Malayalam for to pluck.
\end{quote}

\subsection*{You'll want to know: \ml{കറക്കുക}}
\textbf{\ml{കറക്കുക}} — \emph{kaṟakkuka} — \textquotedblleft{}to milk (a cow)\textquotedblright{}.

\begin{grammarlens}[title={What you've built}]
One more word for this chapter.
\end{grammarlens}

\subsection*{Your turn}
\begin{itemize}
\raggedright
  \item \emph{Say it:} \emph{kaṟakkuka}
  \item \emph{Say it:} \emph{kaṟakkuka}, once more
  \item \emph{Say it:} say \emph{kaṟakkuka}
  \item \emph{Recall:} say the Malayalam for to bend down, then the Malayalam for to pluck, then say \emph{kaṟakkuka} again
\end{itemize}

\begin{tcolorbox}[breakable,colback=teal!4,colframe=teal!35!black,title={Before you move on}]
What does \ml{കറക്കുക} mean? (\textquotedblleft{}To milk (a cow)\textquotedblright{}.) Say it once more.
\end{tcolorbox}
\section[(dialogue)]{Words from chapters 223-226 — a first pass}

\label{lesson:ML-R231-first-pass-words-from-chapters-223-226}



\begin{quote}
Say the words for to pluck and to milk.
\end{quote}

\subsection*{Your turn}
Say these aloud:

\begin{itemize}
\raggedright
  \item to carry, to put, to wear, to take off, to change
  \item to mix, to serve, to taste, to fill, to cover
  \item to drive, to move, to arrive, to return, to take
  \item to discuss, to explain, to prepare, to finish, to check
\end{itemize}

\begin{tcolorbox}[breakable,colback=teal!4,colframe=teal!35!black,title={Before you move on}]
To carry? (\textbf{\ml{ചുമക്കുക}}, \emph{cumakkuka}.) To put? (\textbf{\ml{വയ്ക്കുക}}, \emph{vaykkuka}.) To wear? (\textbf{\ml{ധരിക്കുക}}, \emph{dharikkuka}.) To take off? (\textbf{\ml{ഊരുക}}, \emph{ūruka}.) To change? (\textbf{\ml{മാറ്റുക}}, \emph{māṟṟuka}.) To mix? (\textbf{\ml{കലർത്തുക}}, \emph{kalarttuka}.) To serve? (\textbf{\ml{വിളമ്പുക}}, \emph{viḷambuka}.) To taste? (\textbf{\ml{രുചിക്കുക}}, \emph{rucikkuka}.) To fill? (\textbf{\ml{നിറയ്ക്കുക}}, \emph{niṟaykkuka}.) To cover? (\textbf{\ml{മൂടുക}}, \emph{mūṭuka}.) To drive? (\textbf{\ml{ഓടിക്കുക}}, \emph{ōṭikkuka}.) To move? (\textbf{\ml{നീങ്ങുക}}, \emph{nīṅṅuka}.) To arrive? (\textbf{\ml{എത്തുക}}, \emph{ettuka}.) To return? (\textbf{\ml{മടങ്ങുക}}, \emph{maṭaṅṅuka}.) To take? (\textbf{\ml{കൊണ്ടുപോകുക}}, \emph{koṇṭupōkuka}.) To discuss? (\textbf{\ml{ചർച്ച} \ml{ചെയ്യുക}}, \emph{carcca ceyyuka}.) To explain? (\textbf{\ml{വിശദീകരിക്കുക}}, \emph{viśadīkarikkuka}.) To prepare? (\textbf{\ml{തയ്യാറാക്കുക}}, \emph{tayyāṟākkuka}.) To finish? (\textbf{\ml{പൂർത്തിയാക്കുക}}, \emph{pūrttiyākkuka}.) To check? (\textbf{\ml{പരിശോധിക്കുക}}, \emph{pariśōdhikkuka}.)
\end{tcolorbox}
\section[(dialogue)]{Words from chapters 227-231 — a second pass}

\label{lesson:ML-R231-second-pass-words-from-chapters-227-231}



\begin{quote}
Say the words for to note down and to be present.
\end{quote}

\subsection*{Your turn}
Say these aloud:

\begin{itemize}
\raggedright
  \item to note down, to be present, to take part, to lose one's way, to park
  \item to complain, to tell a lie, to promise, to enjoy, to hate
  \item to set, to shine, to wilt, to doubt, to agree
  \item to suggest, to criticise, to imagine, to solve, to be bored
  \item to lean, to knock, to bend down, to pluck, to milk
\end{itemize}

\begin{tcolorbox}[breakable,colback=teal!4,colframe=teal!35!black,title={Before you move on}]
To note down? (\textbf{\ml{കുറിച്ചെടുക്കുക}}, \emph{kuṟicceṭukkuka}.) To be present? (\textbf{\ml{ഹാജരാകുക}}, \emph{hājarākuka}.) To take part? (\textbf{\ml{പങ്കെടുക്കുക}}, \emph{paṅkeṭukkuka}.) To lose one's way? (\textbf{\ml{വഴിതെറ്റുക}}, \emph{vaḻiteṟṟuka}.) To park? (\textbf{\ml{പാർക്ക്} \ml{ചെയ്യുക}}, \emph{pārkkŭ ceyyuka}.) To complain? (\textbf{\ml{പരാതിപ്പെടുക}}, \emph{parātippeṭuka}.) To tell a lie? (\textbf{\ml{നുണ} \ml{പറയുക}}, \emph{nuṇa paṟayuka}.) To promise? (\textbf{\ml{വാഗ്ദാനം} \ml{ചെയ്യുക}}, \emph{vāgdānaṁ ceyyuka}.) To enjoy? (\textbf{\ml{ആസ്വദിക്കുക}}, \emph{āsvadikkuka}.) To hate? (\textbf{\ml{വെറുക്കുക}}, \emph{veṟukkuka}.) To set? (\textbf{\ml{അസ്തമിക്കുക}}, \emph{astamikkuka}.) To shine? (\textbf{\ml{തിളങ്ങുക}}, \emph{tiḷaṅṅuka}.) To wilt? (\textbf{\ml{വാടുക}}, \emph{vāṭuka}.) To doubt? (\textbf{\ml{സംശയിക്കുക}}, \emph{saṁśayikkuka}.) To agree? (\textbf{\ml{യോജിക്കുക}}, \emph{yōjikkuka}.) To suggest? (\textbf{\ml{നിർദ്ദേശിക്കുക}}, \emph{nirddēśikkuka}.) To criticise? (\textbf{\ml{വിമർശിക്കുക}}, \emph{vimarśikkuka}.) To imagine? (\textbf{\ml{സങ്കൽപ്പിക്കുക}}, \emph{saṅkalppikkuka}.) To solve? (\textbf{\ml{പരിഹരിക്കുക}}, \emph{pariharikkuka}.) To be bored? (\textbf{\ml{ബോറടിക്കുക}}, \emph{bōraṭikkuka}.) To lean? (\textbf{\ml{ചാരുക}}, \emph{cāruka}.) To knock? (\textbf{\ml{മുട്ടുക}}, \emph{muṭṭuka}.) To bend down? (\textbf{\ml{കുനിയുക}}, \emph{kuniyuka}.) To pluck? (\textbf{\ml{പറിക്കുക}}, \emph{paṟikkuka}.) To milk? (\textbf{\ml{കറക്കുക}}, \emph{kaṟakkuka}.)
\end{tcolorbox}
